% The output-first catalogue: one drawn specimen per construct family that
% the tutorial and the recipes do not reach.  Every picture here is modelled
% on a benchmark case or a blueprint figure, so the reader learns the idiom
% the corpus actually writes rather than a toy.
\section{Every construct, drawn}\label{ch:catalogue}

The generated tables in Chapter~\ref{ch:reference} name every environment,
command, and key.  Naming is not showing.  This chapter draws one specimen
per family the tutorial and the recipes leave unpictured, in the order a
source is read: the frame first, then the atoms it places, then the
connections between them, then the annotations over them, and last the
setup that extends the alphabet.

Each specimen is modelled on a benchmark case or a blueprint figure, so
what you read is the spelling the corpus writes.  Where a family has words
the corpus does not write, the entry says so rather than inventing a use.

\subsection{Frames other than the flat row}

A picture's frame is its coordinate system.  \tnkey{frame=}\tnval{flat} is
the default and the whole tutorial rides it; \tnval{plane} projects a
sheet.  The third word is \tnval{circle}, which seats the columns of one
row at equal angles on a ring and turns each glyph by its station's
tangent.  The object it draws is a network whose stations stand on a
cycle, and the contraction that closes the cycle is written rather than
seated.

\begin{tnexample}[
  label={tnex:cat-circle},
  formula={\psi_{ijk}=\sum_{\alpha\beta\gamma}
    A^{i}_{\gamma\alpha}B^{j}_{\alpha\beta}C^{k}_{\beta\gamma}},
  caption={three tensors contracted around a cycle},
  index={Each station carries one physical spoke, and the three virtual
    bonds close the triangle.  The signature records each spoke at the
    bearing its station transported it to, not at the face it was
    declared on.},
  signature={kernel-boundary|signature=phys:210, phys:330, phys:n}]
% Formula: psi_ijk = sum A^i_{gamma alpha} B^j_{alpha beta}
%   C^k_{beta gamma}.
% Ink: the circle frame seats the three sites and carries each
%   physical port out along its station's radius.
\begin{tenkz}[rows={wire}, cols=3, frame=circle, bonds=none]
  \tn[at=(1,1), name=a,
      ports={0:virtual, 180:virtual, 90:physical:$i$}]{}
  \tn[at=(1,2), name=b,
      ports={0:virtual, 180:virtual, 90:physical:$j$}]{}
  \tn[at=(1,3), name=c,
      ports={0:virtual, 180:virtual, 90:physical:$k$}]{}
  \tnwire[name=al]{a.0}{b.180}
  \tnwire[name=be]{b.0}{c.180}
  \tnwire[name=ga]{c.0}{a.180}
  \tnmark[form=label, label pos=e]{on al 0.5}{$\alpha$}
  \tnmark[form=label, label pos=s]{on be 0.5}{$\beta$}
  \tnmark[form=label, label pos=w]{on ga 0.5}{$\gamma$}
\end{tenkz}
\end{tnexample}

Three facts of the ring are visible at once.  A station's local faces are
its own: \tnval{0} is the direction the glyph at that station calls east,
not the page's east, so the one wire \tnval{a.0} to \tnval{b.180} means
the same thing at every station and the three bonds are written alike.
The physical port declared at \tnval{90} leaves along the station's
outward radius, because the row's north is the ring's outward normal.  And
the names read level, because a glyph turns and its label does not.

\tnkey{bonds=}\tnval{none} appears here for the reason it appears in
almost every ring in the corpus: the frame's adjacent-cell contraction
would join the columns in a line, which on a ring is one bond short of the
cycle.  Write the cycle's edges and the picture says what it means.

\dbendpar A ring whose closure is the trace itself, with no tensor on the
return, is common enough to have a word.  \tnkey{ring=}\meta{n} is sugar
for \tnkey{rows=}\tnval{\{wire\}}, \tnkey{cols=}\meta{n},
\tnkey{frame=}\tnval{circle}, \tnkey{west=}\tnval{trace},
\tnkey{east=}\tnval{trace}: the stations of a circle, closed to
themselves.

\begin{tnexample}[
  formula={\operatorname{tr}(M_1M_2M_3M_4)},
  caption={the traced word, drawn as a ring},
  index={The trace follows the frame's arc rather than the chord between
    the two stations it joins.  No index survives, so the picture is a
    scalar.  The names are inscribed: a label outside a glyph on a ring
    is placed toward the centre, where four of them would meet.},
  signature={kernel-boundary|signature=}]
% Formula: tr(M_1 M_2 M_3 M_4); the ring word is sugar for
%   rows={wire}, cols=4, frame=circle, west=trace, east=trace.
\begin{tenkz}[ring=4]
  \tn[skin=box]{M_1} & \tn[skin=box]{M_2} &
  \tn[skin=box]{M_3} & \tn[skin=box]{M_4}
\end{tenkz}
\end{tnexample}

The remaining frame spelling is the explicit one.  \tnkey{basis=} lists
ordered row-kind members at quarter-pitch offsets from the origin, and
every cell address then gains a member coordinate: \tnval{(r,c,k)} is the
$k$-th member's copy of cell \tnval{(r,c)}.  The \tnval{planes} word of
the tutorial's double layer is one two-member basis under a name; writing
the basis out places the members where the mathematics wants them.

\begin{tnexample}[
  caption={three basis members in one frame},
  index={One cell, three members.  The offsets are quarter-pitch steps
    from the origin member, and each member's copy of the cell is
    addressed by its index.},
  signature={kernel-boundary|signature=}]
% Ink: an explicit three-member basis; the member index is the
%   third coordinate of every address.
\begin{tenkz}[rows={wire}, cols=1, bonds=none,
  frame={flat, basis={wire at (0,0), wire at (2,0),
                      wire at (1,2)}}]
  \tn[at=(1,1,1), skin=box]{A}
  \tn[at=(1,1,2), skin=box]{B}
  \tn[at=(1,1,3), skin=box]{C}
\end{tenkz}
\end{tnexample}

\dbendpar An explicit basis belongs to a picture-level \tnval{flat} or
\tnval{plane} frame.  A circle frame refuses one, because its members
would need a tangent-offset contract the language does not state, and a
group frame refuses one because a basis is picture-scoped.  An atom has
no frame of its own at all: neither \tnkey{frame=} nor \tnkey{basis=} is
an atom key, so a site that writes either is turned away as an unknown
word rather than accepted and restricted.

\subsection{The atom record}

An atom's silhouette and its type are separate data.  \tnkey{skin=} names
the silhouette alone; the stock vocabulary is eight words, and the
standard prelude declares two more.

\begin{tnexample}[
  caption={the stock silhouettes, and the two declared by the prelude},
  index={Left to right: dot, box, roundrect, ring, tri, triwest, dots,
    none, and the prelude's pill and mpo.  The eighth cell is empty
    because \tnval{none} draws neither glyph nor label.  None of the ten
    says anything about what a record contracts; that is a declared
    skin's \tnkey{pairings=}, which takes ports and appears with the
    declarations at the end of the chapter.},
  signature={kernel-boundary|signature=}]
% Ink: one atom per skin; bonds=none, so nothing here contracts.
% The none cell is deliberately blank -- that is what it draws.
\begin{tenkz}[rows={wire}, cols=10, bonds=none]
  \tn[skin=dot]{} & \tn[skin=box]{B} & \tn[skin=roundrect]{R} &
  \tn[skin=ring]{X} & \tn[skin=tri]{T} & \tn[skin=triwest]{W} &
  \tn[skin=dots]{} & \tn[skin=none]{} & \tn[skin=pill]{P} &
  \tn[skin=mpo]{O}
\end{tenkz}
\end{tnexample}

Two of these are structural rather than decorative.  \tnval{dots} is the
ellipsis cell of a chain whose middle is elided, and its wires pass
through it; \tnval{none} draws no glyph at all, which is what a corner
junction wants when the wire must turn at a place that carries no tensor.
Both appear in the traced word below.

A glyph's extent is \tnkey{size=}, and a record that stands for a block of
records is \tnkey{cluster=}: an $R\times C$ group of addressable
sub-atoms named after the host.

\begin{tnexample}[
  caption={three size classes, and one clustered site},
  index={The first three are one skin at the three metric classes.  The
    fourth is a single declaration standing for four addressable
    sub-atoms, reached as \tnval{q-1-1} through \tnval{q-2-2}.},
  signature={kernel-boundary|signature=}]
% Ink: size= chooses a glyph's extent from the metric table;
%   cluster= expands one declaration into an addressable group.
\begin{tenkz}[rows={wire}, cols=4, bonds=none]
  \tn[at=(1,1), skin=dot, size=s]{} &
  \tn[at=(1,2), skin=dot, size=m]{} &
  \tn[at=(1,3), skin=dot, size=l]{} &
  \tn[at=(1,4), name=q, cluster=2x2]{}
\end{tenkz}
\end{tnexample}

A site that is not there is \tnkey{void=}.  The two words differ in how
much they take away.  \tnval{sealed} removes the site, the bonds that
reached it, and the leg the row's physical policy would have given it.
\tnval{open} removes the glyph alone: the bonds still meet at the station
and the leg still leaves it, so the picture loses a tensor's ink and not
one entry of its boundary.  The glyph goes because a void site's
silhouette falls to \tnval{none}, and that is a default rather than an
override: a site that names its own \tnkey{skin=} draws it.

\begin{tnexample}[
  formula={\lvert\psi\rangle\ \text{with one site removed}},
  caption={a sealed site in a sheet},
  index={The centre site and its four bonds are gone, and the sheet's
    remaining sites keep their physical legs.  A sealed hole exposes
    nothing, so the signature loses exactly the removed site's leg.},
  signature={kernel-boundary|signature=phys:n, phys:n, phys:n, phys:n,
    phys:n, phys:n, phys:n, phys:n}]
% Ink: void=sealed removes the site and the bonds that reached
%   it; void=open would keep both and drop only the glyph.
\begin{tenkz}[lattice={3x3}, frame=plane, physical=up]
  \tn{} & \tn{} & \tn{} \\
  \tn{} & \tn[void=sealed]{} & \tn{} \\
  \tn{} & \tn{} & \tn{}
\end{tenkz}
\end{tnexample}

\tnexercise{ex:cat-void} The eight physical legs above are the eight
surviving sites.  How many entries does the signature carry when the
centre is written \tnkey{void=}\tnval{open} instead, and which of them
are virtual?

\subsection{Connections}

Every connection is a \tncmd{tnwire}.  Its \tnkey{kind=} says which kind:
\tnval{index} is a contraction, the default, and \tncmd{tnbond} is its
one-word spelling; \tnval{string} travels over the picture carrying an
operator index; \tnval{pairing} is a declared skin's own internal route,
and is never authored at a call site.

The path a wire takes is \tnkey{route=}.  \tnval{straight} is the default
and joins the two ends directly.  \tnval{arc} is the word for a return
that joins two perpendicular faces, and it is what the corpus writes to
close a chain below itself.

\begin{tnexample}[
  wide,
  label={tnex:cat-arc},
  formula={\lvert\psi(X)\rangle=\sum_{s_1\ldots s_N}
    \operatorname{tr}\!\left(XA^{s_1}\cdots A^{s_N}\right)\lvert s\rangle},
  caption={a traced word returning below the row},
  index={The two corner junctions carry no tensor, and the elision cell
    passes its wires through.  Every virtual index closes through $X$;
    each drawn site keeps its physical leg.},
  signature={kernel-boundary|signature=phys:n, phys:n, phys:n}]
% Formula: |psi(X)> = sum tr(X A^{s_1} ... A^{s_N}) |s>.
% Ink: the word closes below the row through two skin=none
%   corner junctions, with X inserted on the return.
\begin{tenkz}[rows={wire,wire}, cols=7, bonds=none]
  \tn[at=(1,2), skin=box, name=a1,
      ports={0:virtual, 180:virtual, 90:physical}]{A}
  \tn[at=(1,3), skin=box, name=a2,
      ports={0:virtual, 180:virtual, 90:physical}]{A}
  \tn[at=(1,5), skin=dots, name=d,
      ports={0:virtual, 180:virtual}]{}
  \tn[at=(1,6), skin=box, name=a4,
      ports={0:virtual, 180:virtual, 90:physical}]{A}
  \tn[at=(2,4), skin=box, name=x,
      ports={0:virtual, 180:virtual}]{X}
  \tn[at=(2,1), skin=none, name=jw,
      ports={0:virtual, 90:virtual}]{}
  \tn[at=(2,7), skin=none, name=je,
      ports={90:virtual, 180:virtual}]{}
  \tnwire{a1.0}{a2.180} \tnwire{a2.0}{d.180}
  \tnwire{d.0}{a4.180}
  \tnwire[route=arc]{a4.0}{je.90}
  \tnwire{je.180}{x.0} \tnwire{x.180}{jw.0}
  \tnwire[route=arc]{jw.90}{a1.180}
\end{tenkz}
\end{tnexample}

Each return joins a pair of perpendicular faces, and the two pairs are
not the same one: the first leaves an east face and enters a north face,
the second leaves a north face and enters a west face.

\dbendpar The word bends the rail exactly where a face directs it.
An arc leaves and enters along its ends' faces, so the two returns
above round their corner junctions instead of cutting across them.  An
end that names no face---a bare name, a coordinate---answers the chord,
and an arc between two such ends is the straight bond it always was.
\tnval{orth} remains the right-angled spelling, drawn from turns
rather than tangents.

The same family draws a closed cycle of tensors laid out as a plaquette,
and a matrix inserted on one of its edges is an ordinary ring atom
addressed by a fraction along the named wire.

\begin{tnexample}[
  formula={\lambda\ \text{on one edge}},
  caption={four arcs, and a bead on one of them},
  index={The four operator-species arcs form the closed ring; the ring
    atom stands halfway along the named north-west arc.  Each tensor
    spends two of its four faces on the ring and leaves two open, so
    eight virtual ends reach the boundary.},
  signature={kernel-boundary|signature=open:e, open:e, open:n, open:n,
    open:s, open:s, open:w, open:w}]
% Ink: four arcs close the plaquette; at=on <wire> <fraction>
%   stands the inserted matrix halfway along one of them.
\tndeclare{species}{op}{hue=source:red}
\begin{tenkz}[rows={wire,wire,wire}, cols=5, bonds=none]
  \tn[skin=mpo, at=(1,3), name=N,
      ports={0:virtual, 90:virtual, 180:virtual, 270:virtual}]{}
  \tn[skin=mpo, at=(2,4), name=E,
      ports={0:virtual, 90:virtual, 180:virtual, 270:virtual}]{}
  \tn[skin=mpo, at=(3,3), name=S,
      ports={0:virtual, 90:virtual, 180:virtual, 270:virtual}]{}
  \tn[skin=mpo, at=(2,2), name=W,
      ports={0:virtual, 90:virtual, 180:virtual, 270:virtual}]{}
  \tnwire[route=arc, species=op, name=lam]{W.90}{N.180}
  \tnwire[route=arc, species=op]{N.0}{E.90}
  \tnwire[route=arc, species=op]{E.270}{S.0}
  \tnwire[route=arc, species=op]{S.180}{W.270}
  \tn[skin=ring, at=on lam 0.5]{\lambda}
\end{tenkz}
\end{tnexample}

The third route word is \tnval{orth}, which turns at right angles.  It is
the spelling for a string that must reach around the picture rather than
across it, and it takes its corners from \tnkey{via=}, a list of addresses
the path passes through.

\begin{tnexample}[
  caption={a string routed around the picture},
  index={The waypoint is an ordinary cell address.  The path leaves $A$
    southward, turns at the named cell, and enters $B$ from the west.},
  signature={kernel-boundary|signature=}]
% Ink: route=orth turns at right angles; via= names the cells
%   the path passes through.
\tndeclare{species}{gauge}{hue=source:blue}
\begin{tenkz}[rows={wire,wire,wire}, cols=4, bonds=none]
  \tn[at=(1,1), name=a, skin=box]{A}
  \tn[at=(3,4), name=b, skin=box]{B}
  \tnwire[kind=string, species=gauge, route=orth,
          via={(3,1)}]{a.270}{b.180}
\end{tenkz}
\end{tnexample}

A string that returns to itself is \tnkey{closed}, and a closed wire takes
no positional ends at all: it is a cycle, so it has none to give.  Its
shape comes from \tnkey{via=} alone.

\begin{tnexample}[
  caption={a closed string around a site},
  index={The loop carries an operator index and contracts nothing.  It
    exposes no end, so it adds no entry to the signature.},
  signature={kernel-boundary|signature=}]
% Ink: a closed string takes no ends; via= gives it its shape.
\begin{tenkz}[rows={wire}, cols=2, bonds=none]
  \tn[at=(1,1), name=z, skin=box]{Z}
  \tnwire[kind=string, closed, via={0.6 n of z, 0.6 e of z,
    0.6 s of z, 0.6 w of z}]
\end{tenkz}
\end{tnexample}

Three further wire keys decorate the rail rather than change its path.
\tnkey{stroke=} is the rail's pattern, and the corpus draws only three:
\tnval{solid} for a contraction, \tnval{dashed} for a line the figure
shows but its contraction does not use, \tnval{dotted} for a lattice
lying under the sheet being drawn.  Those are readings, and the pattern
is ink alone: a dashed wire keeps the default
\tnkey{kind=}\tnval{index}, and what decides whether a drawn line
contracts is what its two ends are.  A wire between two ports takes them
both, whatever its stroke; a wire between two atoms named as wholes
leaves their ports where they were.  That is how the dashed diagonals of
the benchmark's dual lattice recall an edge without contracting one:
their ends are junctions with no ports to take.  \tnkey{weight=} says how
many indices one drawn rail stands for, and \tnkey{dir=} marks a directed
index against its dual.

\begin{tnexample}[
  wide,
  caption={the three strokes, a bundle, and a direction},
  index={Left: the three declared rail patterns.  The four sites declare
    no ports, so the three rails differ in pattern and in nothing else.
    Right: one rail standing for three indices, carrying the direction
    mark that distinguishes a space from its dual.  The bundle is
    recorded, not drawn: the event stream reads \tnval{weight=bundle=3}
    and the page shows one line.},
  signature={kernel-boundary|signature=}]
% Ink: stroke= is the rail's pattern; weight=bundle=n records
%   that one rail carries n indices; dir= marks the direction.
\begin{tenkz}[rows={wire}, cols=4, bonds=none]
  \tn[at=(1,1), name=p]{} \tn[at=(1,2), name=q]{}
  \tn[at=(1,3), name=r]{} \tn[at=(1,4), name=s]{}
  \tnwire[stroke=solid]{p.0}{q.180}
  \tnwire[stroke=dashed]{q.0}{r.180}
  \tnwire[stroke=dotted]{r.0}{s.180}
\end{tenkz}
\qquad
\begin{tenkz}[rows={wire}, cols=2, bonds=none]
  \tn[at=(1,1), name=u]{U} \tn[at=(1,2), name=v]{V}
  \tnwire[weight=bundle=3, dir=to]{u.0}{v.180}
\end{tenkz}
\end{tnexample}

\subsection{Annotations}

A mark speaks over a picture and owns none of it.  Its \tnkey{form=} is
the shape of the speech, and the alphabet is four words.  Three put ink
on the page: \tnval{label} stands text at a target, \tnval{bracket}
measures a span from outside it, and \tnval{enclosure} strokes the offset
hull of a selection.  The fourth, \tnval{prose}, is accepted and recorded
without ink.  This section draws the three, shows what a spelling the
alphabet has lost answers now, and reads \tnval{prose} out of the record
stream.

\begin{tnexample}[
  caption={a bracket over a span, a label on a site},
  index={The bracket names the range of two cells; the label names one.
    Neither touches the contraction, and the signature is the row's.},
  signature={kernel-boundary|signature=open:e, open:w, phys:n, phys:n,
    phys:n}]
% Ink: two mark forms over one row; no contraction changes.
\begin{tenkz}[rows={wire}, cols=3, west=open, east=open,
              physical=up]
  \tn{A} & \tn{B} & \tn{C}
  \tnmark[form=bracket]{(1,1) .. (1,2)}{$P$}
  \tnmark[form=label, label pos=n]{(1,3)}{$q$}
\end{tenkz}
\end{tnexample}

Contours nest by themselves.  A second enclosure over one selection
stands one clearance outside the first because it was declared second,
which is how a figure draws a boundary operator acting on the open
indices of a region it also has to show.

\begin{tnexample}[
  formula={\lvert\psi(X)\rangle},
  caption={a region, and an operator on its boundary},
  index={The inner contour is the retained block; the outer one, declared
    second and therefore one clearance further out, is the operator
    acting on the indices the block leaves open.},
  signature={kernel-boundary|signature=phys:n, phys:n, phys:n, phys:n}]
% Ink: two contours over one selection; concentric order steps
%   the second declared clear of the first.
\begin{tenkz}[rows={wire}, cols=4, bonds=none, physical=up]
  \tn[at=(1,1), skin=box]{A} & \tn[at=(1,2), skin=box]{A} &
  \tn[at=(1,3), skin=box]{A} & \tn[at=(1,4), skin=box]{A}
  \tnmark[form=enclosure]{(1,1) .. (1,4)}{}
  \tnmark[form=enclosure, label pos=e]{(1,1) .. (1,4)}{$X$}
\end{tenkz}
\end{tnexample}

Nothing in the source states that order.  Contours are ranked by
containment, ties going to the earlier declaration, and each rank stands
one clearance further out than the rank inside it.  A bracket drawn
inside a region window is ranked the same way and clears it instead of
sharing its support.

\dbendpar The registry once carried an \tnkey{inset=} key for this,
named as the nested-contour step.  It recorded a number in the event
stream and moved no ink -- write one, write six, or leave the key out,
and the two contours above fell in the same places -- so the word is
retired: a document that writes it meets the unknown-key refusal, and
the tombstone ledger names what the sources get instead, which is the
ranking above and nothing to set.

\tnkey{slot=} tints the contour of an enclosure or a bracket from the
theme's semantic slots, and reaches no other form: a label takes its ink
from \tnkey{species=} alone, so a slot written on a label is accepted,
recorded, and drawn in the ordinary ink.  \tnkey{tint} lays ink over the
paper inside a contour rather than only stroking it; the tutorial's
blocked run shows the tinted form, and so does the band below.

Four further words stood in the \tnkey{form=} list until the alphabet
closed on the four above.  Each of them took a record and drew nothing,
and each is now a retired spelling that names its successor.  A brace is
a \tnval{bracket}, which speaks from the side its label sits on, so
\tnval{brace-above} is that bracket named to the north and
\tnval{brace-below} is the same bracket with no side named, the south
being where a bracket speaks by default.  A \tnval{band} is an
\tnval{enclosure} carrying \tnkey{tint}, the paper under the contour
being the whole of what a band added to one.  A \tnval{cut} has no
successor of its own: a cut in the sources is the contour of what it
separates, which is an enclosure over that selection.  The two brace
spellings and the band are drawn below in the words that replaced them.

\begin{tnexample}[
  caption={a brace above, a brace below, and a band},
  index={One bracket named on two sides, and one tinted enclosure.  A
    bracket's arc follows its name, so \tnkey{label pos=}\tnval{90}
    carries both to the north and the default carries both to the south.
    The enclosure's \tnkey{tint} is the paper a band laid under its
    span.},
  signature={kernel-boundary|signature=}]
% Ink: the two brace spellings are one bracket named on two
%   sides; the band is an enclosure carrying tint.
\begin{tenkz}[rows={wire}, cols=2]
  \tn[skin=box]{A} & \tn[skin=box]{A}
  \tnmark[form=bracket, label pos=90]{(1,1) .. (1,2)}{$N$}
\end{tenkz}
\qquad
\begin{tenkz}[rows={wire}, cols=2]
  \tn[skin=box]{A} & \tn[skin=box]{A}
  \tnmark[form=bracket]{(1,1) .. (1,2)}{$N$}
\end{tenkz}
\qquad
\begin{tenkz}[rows={wire}, cols=2]
  \tn[skin=box]{A} & \tn[skin=box]{A}
  \tnmark[form=enclosure, tint]{(1,1) .. (1,2)}{}
\end{tenkz}
\end{tnexample}

A source that still writes one of the four is refused by the word it
wrote, and the answer carries the migration rather than the alphabet:

\begin{tnrefusal}[
  caption={a retired form, refused with its successor},
  index={The parser still knows a struck word, so that it can answer
    for it.  Reporting the alphabet alone would list the words that
    remain, which reads as a typo when the word was in fact removed.
    Each of the four answers with its own migration, and \tnval{cut}'s
    says it has no successor.},
  diagnostic={[TKZ-LANG-TOMBSTONE] form=brace-above retired from the
    alphabet: use form=bracket with label pos=90; a bracket speaks from
    the side its label sits on (record mark-4).}]
% Refused: brace-above is a retired spelling.  A brace above a
% span is form=bracket with label pos=90.
\begin{tenkz}[rows={wire}, cols=3, west=open, east=open]
  \tn{A} & \tn{B} & \tn{C}
  \tnmark[form=brace-above]{(1,1) .. (1,2)}{$P$}
\end{tenkz}
\end{tnrefusal}

The migration is written once, as a tombstone row of the registry, and
the parser answers in that row's words, so a reader of an older source
meets the same sentence wherever it is read.  It is also why a retired
word is refused rather than accepted and drawn as nothing: a picture
either carries the mark its source asked for or does not compile.

\tnval{prose} is the one form that is accepted without ink, and it is
accepted because it asks for none.  \tncmd{tnprose} states, in the
picture's own record stream, that a step of an argument was not drawn.
Nothing appears on the page and one line appears in the event stream:

\begin{center}
\tnlogline{mark|id=mark-1|form=prose|label=by~isometry|target=panel}
\end{center}

\noindent On a checked relation that line is a hard error without a
recorded opt-out, because prose is not a boundary signature and cannot be
compared with one.

\subsection{Fusion}

Two atoms are structured rather than atomic.  \tncmd{tntree} takes a
parenthesized fusion word and draws the tree that parses it; the strands
are drawn as lines by default, and \tnkey{tree style=}\tnval{ribbon}
draws the same parsed tree as a regular ribbon neighbourhood, one band per
simple object.

\begin{tnexample}[
  formula={m_{ab}^{x}\colon a\otimes b\to x},
  caption={one fusion tree in both strand styles},
  index={Both pictures are the same chosen fusion morphism.  The wire
    style meets three lines at one trivalent vertex; the ribbon style
    gives each simple object a band and joins them at a vertex patch.},
  signature={}]
% Ink: one parsed fusion word, drawn in both strand styles.
\tndeclare{species}{anyon}{}
$\tntree[tree style=wire, species=anyon]{(a\,b)_x}
 \qquad
 \tntree[tree style=ribbon, species=anyon]{(a\,b)_x}$
\end{tnexample}

A tree is a self-contained box and stands wherever mathematics stands: in
a display, in an alignment cell, in a \texttt{tikz-cd} cell.  It is the
one construct here that does not need a picture around it.

\tncmd{tnfuse} is the other.  It is a fusion atom, a triangle whose split
side the frame bonds to the rows it spans, and it is how a doubled-layer
object is written as one tensor.

\begin{tnexample}[
  formula={M^{ij}=\sum_{k}e\,(A^{(i,k)}\otimes\overline{A}^{(j,k)})
    \,e^{-1}},
  caption={a fuser, the doubled layer, and its inverse},
  index={The two fusion atoms span both rows and combine the ket and bra
    virtual indices into one; between them the two layers carry the
    physical indices $i$ and $j$.  The frame's bond between the rows is
    the ancillary index $k$ the sum runs over.},
  signature={kernel-boundary|signature=open:e, open:w, phys:n, phys:s}]
% Formula: M^{ij} = sum_k e (A^{(i,k)} tensor
%   conj(A)^{(j,k)}) e^{-1}; the fusers combine the ket and bra
%   virtual indices into one, and the bond between the rows is
%   the ancillary index k.
\begin{tenkz}[rows={op,op}, cols=3]
  \tnfuse[at=(1,1), name=e, skin=pill]{e}
  \tn[at=(1,2), skin=box, ports={90:physical:$i$}]{A}
  \tn[at=(2,2), skin=box, ports={270:physical:$j$}]{\overline{A}}
  \tnfuse[at=(1,3), name=ei, skin=pill]{e^{-1}}
  \tnwire{open w}{e.180} \tnwire{ei.0}{open e}
\end{tenkz}
\end{tnexample}

\dbendpar \tncmd{tnfuse} is sugar for
\tnval{\textbackslash tn[skin=tri, wires=2, \ldots]}, and authored keys
follow the preset, so the \tnkey{skin=}\tnval{pill} above overrides the
triangle the preset would have drawn.  Under
\tncmd{tnset}\tnval{\{strict\}} the sugar is refused and the kernel
spelling is required.

\subsection{Setup and extension}

\tncmd{tndeclare} is the one extension door, and it has three classes.
A \emph{species} is a named hue a paper's family of objects keeps across
figures.  A \emph{skin} is a base silhouette plus \tnkey{pairings=}, its
own internal routes between its own ports.  An \emph{atom} is a new typed
spelling in the atom class.

The skin class is the one that changes what a glyph means.  A tensor whose
two port pairs are joined inside the box is drawn once as a declaration
and used as a word:

\begin{tnexample}[
  wide,
  formula={U=U_{\rightarrow}\,U_{\leftarrow}},
  caption={two declared skins, each with its own internal routes},
  index={The upper row turns its north input east and its south input
    west; the lower row reverses both.  The pairings belong to the skin,
    so every site drawn with it carries them.  The row policy grows a leg
    on both faces of every site, and the frame's bond between the rows
    takes the inner one at each, so six legs reach the boundary and not
    twelve.},
  signature={kernel-boundary|signature=phys:n, phys:n, phys:n, phys:s,
    phys:s, phys:s}]
% Ink: a declared skin is a base silhouette plus its own
%   internal port pairings, tinted by declared species.
\tndeclare{species}{right}{hue=source:blue}
\tndeclare{species}{left}{hue=source:red}
\tndeclare{skin}{shift-right}{
  base=box, pairings={n@1 > e@1 : right, s@1 > w@1 : left}}
\tndeclare{skin}{shift-left}{
  base=box, pairings={n@1 > w@1 : left, s@1 > e@1 : right}}
\begin{tenkz}[rows={wire,wire}, cols=3, physical=updown]
  \tn[skin=shift-right]{} & \tn[skin=shift-right]{} &
  \tn[skin=shift-right]{} \\
  \tn[skin=shift-left]{} & \tn[skin=shift-left]{} &
  \tn[skin=shift-left]{}
\end{tenkz}
\end{tnexample}

A pairing is a wire whose two ends are the skin's own ports, so it is
reusable topology: declare it once and every site drawn with that skin
carries it.  \tnkey{base=} may name another declared skin, so a family of
related glyphs is written as a chain of small declarations rather than
one long one.  What the chain carries down is the silhouette alone: a
derived skin resolves its base to a primitive and inherits none of its
base's pairings, so a route the derived glyph wants is written again in
its own declaration.

The atom class has a two-argument spelling, \tncmd{tndeclareatom}, for the
common case of a new command with a skin and a port list.

\begin{tnexample}[
  caption={a declared atom command},
  index={The declaration adds one word to the atom class; the picture
    then spells it like any other atom.  The port list is complete: an
    anchor without a type is not an extension contract.},
  signature={kernel-boundary|signature=open:e, open:w, phys:n}]
% Ink: tndeclareatom names faces west|east:virtual and
%   up|down:physical; tndeclare{atom} also accepts angles.
\tndeclareatom{\tnprojector}{skin=box,
  ports={west:virtual, east:virtual, up:physical}}
\begin{tenkz}[rows={wire}, cols=3, west=open, east=open]
  \tn{A} & \tnprojector{P} & \tn{B}
\end{tenkz}
\end{tnexample}

\dbendpar The two spellings do not accept the same vocabulary, in ports
or in skins.  \tncmd{tndeclareatom} takes the four named faces only ---
\tnval{west}, \tnval{east} typed \tnval{virtual}, and \tnval{up},
\tnval{down} typed \tnval{physical} --- and answers anything else with
\tnval{[TKZ-ATOM-INVALID-PORT]}.  Its \tnkey{skin=} is four words as
well, \tnval{dot}, \tnval{box}, \tnval{pill} and \tnval{mpo}; a fifth is
refused as an unknown choice and the declaration falls back to
\tnval{dot}.  The three-argument \tncmd{tndeclare}\tnval{\{atom\}} also
takes compass letters, numeric angles, and every silhouette, so a glyph
with an off-axis face or a \tnval{ring} skin is declared there.

A declaration is setup, and setup happens outside a composition form.
Writing one among an equation's panels is refused: it would be a
document-wide act performed in the middle of an assertion, ordered by
where it happens to sit and read once for each time the equation is
measured.  Write it before the equation, where it says the same thing
once.

\tncmd{tnset} carries the same policy at document scope.  Four keys live
there: \tnkey{pitch=}, the one base metric every named ratio follows;
\tnkey{sizes=}, its size-class table; \tnkey{theme=}, the semantic ink;
and \tnkey{strict}, which rejects every sugar spelling in the registry.
No benchmark case sets it.  The canonical cases write sugar freely, and
\tnkey{physical=} alone stands in fifty-one of the hundred and thirty of
them; the sources that turn \tnval{strict} on are the five refusal
fixtures that prove it refuses, together with the refusal below, which is
a sixth of the same kind.
What the checks hold is the sugar itself: eleven sugar spellings compile
beside their kernel expansions and must emit the same events and the same
picture, so a sugar word is a shorter way to write a canonical case and
never a second meaning.

\begin{tnrefusal}[
  caption={strict setup rejects a sugar spelling},
  index={\tnkey{sandwich} is sugar for
    \tnkey{rows=}\tnval{\{ket,op,bra\}}.  Under \tnval{strict} the
    expansion is not performed; the spelling is refused and named.},
  diagnostic={[TKZ-LANG-STRICT] 'sandwich' is sugar and
    \textbackslash tnset\{strict\} is active.}]
% Refused: strict rejects sugar.  The kernel spelling of this
% picture is rows={ket,op,bra}.
\tnset{strict}
\begin{tenkz}[sandwich, cols=1]
  \tn{A} \\ \tn[skin=mpo]{O} \\ \tn{$\overline{A}$}
\end{tenkz}
\end{tnrefusal}

One composition word closes the alphabet.  \tncmd{tngroup} transforms a
sub-diagram as one object, and the records inside it keep the names they
were given, so the rest of the picture can still address them.

\begin{tnexample}[
  caption={a grouped sub-diagram},
  index={The group is one object of its picture.  Its member keeps the
    name \tnval{g} and is reachable from outside the group by that name.},
  signature={kernel-boundary|signature=}]
% Ink: a group is a sub-frame of its picture, and the names
%   declared inside it stay addressable outside it.
\begin{tenkz}[rows={ket}, cols=1, bonds=none]
  \tngroup[frame=flat]{ \tn[skin=box, name=g]{G} }
\end{tenkz}
\end{tnexample}

A group's options are picture keys, and \tnkey{frame=} is the one the
specimen writes, a sub-frame being what a group is for.  A group shares
its picture's metric, size class, and audit, so the three keys that would
set those over again are refused at group scope rather than accepted and
left inert: \tnkey{size=}, \tnkey{metrics=}, and \tnkey{check=}.  A
\tnkey{basis=} inside a group's frame is refused with them.

\subsection{Answers to the exercises}

\tnanswer{ex:cat-void} Nine entries, none of them virtual, which is the
signature of the sheet with no hole in it at all.  \tnval{open} takes
away the glyph and nothing else: the four bonds still meet at the empty
station, and the centre still carries the leg its row's policy gives it,
so a reader who wants the eight-leg boundary above must write
\tnval{sealed}.  Neither word turns a bond into an exposed index.  What
\tnval{open} buys is the picture: a station where the sheet contracts and
no tensor is drawn.
